\documentclass[a4paper]{article}

\usepackage{amsfonts, amsmath, amssymb, amsthm}
\usepackage{fullpage}
\usepackage{hyperref}

\newtheorem{theorem}{Theorem}[section]
\newtheorem{lemma}[theorem]{Lemma}
\newtheorem{corollary}[theorem]{Corollary}
\newtheorem{proposition}[theorem]{Proposition}
\newtheorem{conjecture}[theorem]{Conjecture}
\newtheorem{remark}[theorem]{Remark}

\newcommand{\F}{\mathbb{F}}
\newcommand{\ord}[2][]{\operatorname{ord}_{#1}\!\left(#2\right)}
\newcommand{\rhoord}[1]{\rho_2\!\left(#1\right)}
\renewcommand{\qedsymbol}{$\blacksquare$}

\setlength{\parindent}{0em}
\setlength{\parskip}{1em}

\begin{document}

\title{Prime-Power Stabilization and the Wieferich Obstruction\\
for \textit{Lights Out} Nullity}
\author{William Boyles}
\date{\today}
\maketitle

\begin{abstract}
	Let $d(n)$ denote the nullity over $\F_2$ of the $n \times n$ bounded \textit{Lights Out} grid.
	An existing argument proves
	\[
		d(p^k-1)=d(p-1)
	\]
	for every non-Wieferich prime $p$.
	We show more generally that, if
	\[
		c_p=v_p\!\left(2^{\ord[p]{2}}-1\right),
	\]
	then $d(p^k-1)$ is constant for all $k\geq c_p$.
	Thus every possible exception is confined to the finite initial plateau
	$p,p^2,\ldots,p^{c_p}$ on which the order of $2$ does not grow.
	Exact polynomial computations show that no exception occurs for either known Wieferich prime, 1093 or 3511.
	Consequently the equality holds for every currently known prime, and in particular for every prime below $2^{64}$.
	The argument does not settle the universal statement for a presently unknown Wieferich prime.
\end{abstract}

\section{Fibonacci-polynomial formulation}

Let $F_n(x)\in\F_2[x]$ be defined by
\[
	F_0(x)=0,\qquad F_1(x)=1,\qquad
	F_n(x)=xF_{n-1}(x)+F_{n-2}(x).
\]
Sutner's polynomial formulation of bounded \textit{Lights Out} gives
\begin{equation}\label{eq:nullity-gcd}
	d(n)=\deg\gcd\!\left(F_{n+1}(x),F_{n+1}(x+1)\right)
\end{equation}
\cite{Sutner2000}.
The required divisibility and root properties of these polynomials are developed in
\cite{HUNZIKER2004465}.

For an odd prime $p$ and $j\geq1$, define
\[
	R_j=\left\{\zeta+\zeta^{-1}:\ord{\zeta}=p^j\right\}
	\subset\overline{\F}_2.
\]
The distinct roots of $F_{p^k}$ are
\begin{equation}\label{eq:root-union}
	\bigcup_{j=1}^k R_j.
\end{equation}
Moreover, $F_{p^k}$ is the square of a square-free polynomial, so each of these roots has multiplicity two.

For an odd integer $M$, let
\[
	\rho_2(M)=\min\{r>0:2^r\equiv\pm1\pmod M\}
\]
be the signed order of $2$ modulo $M$.
If $\zeta$ has exact order $M$ and $\alpha=\zeta+\zeta^{-1}$, then
\begin{equation}\label{eq:field-degree}
	[\F_2(\alpha):\F_2]=\rho_2(M).
\end{equation}
Indeed, $\alpha^{2^r}=\alpha$ precisely when $\zeta^{2^r}=\zeta^{\pm1}$.

\section{The stabilization theorem}

We first record the standard order formula that identifies the exceptional plateau.

\begin{lemma}\label{lem:orders}
	Let $p$ be an odd prime, set $h=\ord[p]{2}$, and let
	\[
		c=v_p(2^h-1).
	\]
	Then, for every $j\geq1$,
	\[
		\ord[p^j]{2}=h\,p^{\max(0,j-c)}
	\]
	and
	\[
		\rho_2(p^j)=\rho_2(p)\,p^{\max(0,j-c)}.
	\]
\end{lemma}

\begin{proof}
	If $2^r\equiv1\pmod{p^j}$, then $h\mid r$, so write $r=ht$.
	The lifting-the-exponent lemma gives
	\[
		v_p(2^{ht}-1)=v_p(2^h-1)+v_p(t)=c+v_p(t).
	\]
	The least possible $t$ is therefore $p^{\max(0,j-c)}$.
	This proves the first formula.

	Since $p$ is odd, multiplication by a power of $p$ does not change parity.
	The ordinary orders $\ord[p^j]{2}$ therefore all have the same parity as $h$.
	Because $(\mathbb{Z}/p^j\mathbb{Z})^\times$ is cyclic, the unique element of order two is $-1$.
	The signed order is consequently the ordinary order when that order is odd and half the ordinary order when it is even, which gives the second formula.
\end{proof}

\begin{theorem}[Prime-power stabilization]\label{thm:stabilization}
	Let $p$ be an odd prime and set
	\[
		c=v_p\!\left(2^{\ord[p]{2}}-1\right).
	\]
	Then, for every $k\geq c$,
	\[
		d(p^k-1)=d(p^c-1).
	\]
\end{theorem}

\begin{proof}
	Fix $k\geq c$.
	Suppose $\alpha$ and $\alpha+1$ are both roots of $F_{p^k}$, with
	\[
		\alpha\in R_j,\qquad \alpha+1\in R_\ell.
	\]
	Because $\F_2(\alpha)=\F_2(\alpha+1)$, their extension degrees over $\F_2$ are equal.
	Equations \eqref{eq:field-degree} and Lemma \ref{lem:orders} give
	\[
		\rho_2(p)\,p^{\max(0,j-c)}
		=
		\rho_2(p)\,p^{\max(0,\ell-c)}.
	\]
	Hence either $j,\ell\leq c$, or $j=\ell>c$.

	It remains to exclude the second possibility.
	Assume at least one of $j,\ell$ is greater than $c$.
	The degree equality above then forces $j=\ell>c$, so in particular
	\[
		\alpha+1\in R_\ell=R_j.
	\]
	Set
	\[
		M=p^j,\qquad m=p^{j-c},
	\]
	and choose a primitive $M$th root $\zeta$ such that
	\[
		\alpha=\zeta+\zeta^{-1}.
	\]
	Since $\alpha+1\in R_j$, there is an integer $s$ with $\gcd(s,p)=1$ such that
	\[
		\alpha+1=\zeta^s+\zeta^{-s}.
	\]
	Thus
	\[
		1+\zeta+\zeta^{-1}+\zeta^s+\zeta^{-s}=0,
	\]
	so $\zeta$ is a root of
	\begin{equation}\label{eq:sparse-polynomial}
		P_s(x)=1+x+x^{M-1}+x^s+x^{M-s}.
	\end{equation}

	The element $\zeta^m$ has exact order $p^c$.
	Let $f(x)$ be its minimal polynomial over $\F_2$.
	By Lemma \ref{lem:orders},
	\[
		\deg f=\ord[p^c]{2}=h.
	\]
	Now $\zeta$ is a root of $f(x^m)$, and
	\[
		\deg f(x^m)=hm
		=h p^{j-c}
		=\ord[p^j]{2}.
	\]
	This is exactly the degree of the minimal polynomial of $\zeta$.
	Consequently $f(x^m)$ is that minimal polynomial, and
	\[
		f(x^m)\mid P_s(x).
	\]

	Every polynomial has a unique decomposition by exponents modulo $m$:
	\[
		P_s(x)=\sum_{a=0}^{m-1}x^aP_a(x^m).
	\]
	If $f(x^m)\mid P_s(x)$, uniqueness implies that $f(y)$ divides every $P_a(y)$.
	However, only the constant exponent in \eqref{eq:sparse-polynomial} is divisible by $m$.
	Indeed, the other four exponents are congruent modulo $m$ to $\pm1$ and $\pm s$, all nonzero because $p\nmid s$.
	Therefore $P_0(y)=1$, which cannot be divisible by the nonconstant polynomial $f(y)$.
	This contradiction excludes every common root in $R_j$ for $j>c$.

	All common roots of $F_{p^k}(x)$ and $F_{p^k}(x+1)$ therefore lie in
	\[
		\bigcup_{j=1}^cR_j,
	\]
	the root set of $F_{p^c}$.
	Conversely, each common root at a level at most $c$ remains a common root at level $k$.
	Every root has multiplicity two in both Fibonacci polynomials, so their GCDs have the same degree.
	Equation \eqref{eq:nullity-gcd} now gives
	\[
		d(p^k-1)=d(p^c-1).
	\]
\end{proof}

\begin{corollary}\label{cor:non-wieferich}
	If $p$ is a non-Wieferich prime, then for every $k\geq1$,
	\[
		d(p^k-1)=d(p-1).
	\]
\end{corollary}

\begin{proof}
	For an odd non-Wieferich prime,
	\[
		v_p\!\left(2^{\ord[p]{2}}-1\right)=1,
	\]
	so Theorem \ref{thm:stabilization} applies with $c=1$.
	The case $p=2$ is the known identity $d(2^k-1)=0$ \cite{Sutner1989}.
\end{proof}

\begin{remark}
	Theorem \ref{thm:stabilization} is strictly stronger than the non-Wieferich result.
	It says that a prime can behave exceptionally only while the multiplicative order of $2$ remains on its initial plateau.
	Even for a higher-order Wieferich prime with $c>2$, no new common roots can first appear after exponent $c$.
\end{remark}

\section{The two known Wieferich primes}

The only known base-$2$ Wieferich primes are 1093 and 3511 \cite{CrandallSearch,PrimeGridWieferich}.
For both, the Wieferich depth is exactly two:
\[
	v_{1093}(2^{364}-1)=2,\qquad
	v_{3511}(2^{1755}-1)=2.
\]
Thus Theorem \ref{thm:stabilization} reduces each infinite family to a single new computation at $p^2$.

The base nullities vanish without a large polynomial calculation.
Blokhuis proved that if an odd prime $p$ is mad for the toroidal version of the game, then
\[
	\rho_2(p)\leq\sqrt p
\]
\cite[Theorem 4.2]{ButtonMadness}.
A nonzero bounded-grid nullity $d(p-1)$ produces such a torus solution.
But
\[
	\rho_2(1093)=182>\sqrt{1093},
	\qquad
	\rho_2(3511)=1755>\sqrt{3511}.
\]
Hence
\[
	d(1093-1)=d(3511-1)=0.
\]

\begin{proposition}[Exact computation]\label{prop:known-computation}
	We have
	\[
		d(1093^2-1)=0
		\qquad\text{and}\qquad
		d(3511^2-1)=0.
	\]
\end{proposition}

\begin{proof}
	For odd $m=2r+1$, the fast-doubling identities give
	\[
		F_m(x)=\left(F_r(x)+F_{r+1}(x)\right)^2.
	\]
	Set
	\[
		\tau_m(x)=F_r(x)+F_{r+1}(x).
	\]
	Then $\tau_m$ is square-free and
	\[
		d(m-1)
		=
		2\deg\gcd\!\left(\tau_m(x),\tau_m(x+1)\right).
	\]

	The exact checker \texttt{code/wieferich\_check.py} constructs both polynomials by fast doubling and computes their GCDs in $\F_2[x]$ using FLINT.
	It returns degree zero for $m=1093^2$ and $m=3511^2$.
	All operations are exact; the calculation uses neither randomization nor floating-point arithmetic.
\end{proof}

\begin{corollary}\label{cor:known-wieferich}
	For $p\in\{1093,3511\}$ and every $k\geq1$,
	\[
		d(p^k-1)=d(p-1)=0.
	\]
\end{corollary}

\begin{proof}
	The case $k=1$ and Proposition \ref{prop:known-computation} handle the two levels on the plateau.
	Theorem \ref{thm:stabilization} handles every $k\geq2$.
\end{proof}

\begin{corollary}\label{cor:below-bound}
	For every prime $p<2^{64}$ and every $k\geq1$,
	\[
		d(p^k-1)=d(p-1).
	\]
\end{corollary}

\begin{proof}
	Searches have shown that the only Wieferich primes below $2^{64}$ are 1093 and 3511 \cite{PrimeGridWieferich}.
	Apply Corollary \ref{cor:known-wieferich} to those two primes and Corollary \ref{cor:non-wieferich} to every other prime in the range.
\end{proof}

\section{What remains unresolved}

The stabilization theorem sharply isolates the universal question.

\begin{proposition}\label{prop:finite-obstruction}
	Let $p$ be an odd prime and
	\[
		c=v_p\!\left(2^{\ord[p]{2}}-1\right).
	\]
	The identity
	\[
		d(p^k-1)=d(p-1)\qquad(k\geq1)
	\]
	holds if and only if it holds for the finitely many exponents $1\leq k\leq c$.
\end{proposition}

\begin{proof}
	One implication is immediate.
	For the other, equality at $k=c$ together with Theorem \ref{thm:stabilization} gives equality for every $k\geq c$.
\end{proof}

For $1\leq j,\ell\leq c$, all elements of $R_j$ and $R_\ell$ have the same extension degree over $\F_2$.
This is exactly why the degree-separation argument cannot distinguish the levels on the plateau.
A new common root there is equivalent to a solution
\[
	1+\zeta+\zeta^{-1}+\eta+\eta^{-1}=0
\]
with $\zeta$ and $\eta$ of $p$-power order at most $p^c$, at least one of them having order greater than $p$.
Equivalently, in the same-level case one obtains the sparse cyclotomic relation
\[
	1+\zeta+\zeta^{-1}+\zeta^s+\zeta^{-s}=0
\]
in characteristic two.

The theorem of Lam and Leung on characteristic-zero vanishing sums would force a five-term sum of roots of $p$-power order to have $p=5$ \cite{LamLeung2000}.
It does not resolve the relation above: reduction modulo two can create additive relations that do not lift to characteristic zero.
Likewise, the general Chebyshev-polynomial and $\sigma$-automata results in
\cite{Sutner2000,HUNZIKER2004465} provide the polynomial framework but do not eliminate the Wieferich plateau.

\begin{conjecture}\label{conj:all-primes}
	For every prime $p$ and every $k\geq1$,
	\[
		d(p^k-1)=d(p-1).
	\]
\end{conjecture}

The results above neither prove nor disprove Conjecture \ref{conj:all-primes}.
They show that any counterexample must involve a presently unknown Wieferich prime and must already occur within its finite initial plateau.

\bibliographystyle{amsalpha}
\bibliography{refs.bib}

\end{document}
